\documentclass[12pt,a4paper]{report}
\usepackage[utf8]{inputenc}
\usepackage[T1]{fontenc}
\usepackage[spanish,es-nodecimaldot]{babel}
\usepackage[a4paper,margin=2.45cm,headheight=26pt]{geometry}
\usepackage{lmodern}
\usepackage{graphicx}
\usepackage{xcolor}
\usepackage{booktabs}
\usepackage{longtable}
\usepackage{listings}
\usepackage{tikz}
\usetikzlibrary{positioning,arrows.meta,calc}
\usepackage{float}
\usepackage{caption}
\usepackage{tcolorbox}
\usepackage{fancyhdr}
\usepackage{enumitem}
\usepackage{hyperref}

\definecolor{techpurple}{HTML}{7135A5}
\definecolor{techpurple2}{HTML}{914BC6}
\definecolor{techlavender}{HTML}{F3EFF8}
\definecolor{techink}{HTML}{25212B}
\definecolor{techpink}{HTML}{D94E91}
\definecolor{techline}{HTML}{DED0EA}
\hypersetup{colorlinks=true,linkcolor=techpurple,urlcolor=techpurple2,citecolor=techpurple}
\pagestyle{fancy}
\fancyhf{}
\lhead{\textcolor{techpurple}{\textbf{TECH-FLOW}}}
\rhead{\textcolor{techpurple}{\leftmark}}
\cfoot{\textcolor{techpurple}{\thepage}}
\renewcommand{\headrulewidth}{0.5pt}
\renewcommand{\headrule}{\hbox to\headwidth{\color{techline}\leaders\hrule height \headrulewidth\hfill}}
\setlength{\parskip}{5pt}
\setlength{\parindent}{0pt}
\setlength{\emergencystretch}{3em}
\setlist{itemsep=3pt,topsep=4pt}
\captionsetup{font=small,labelfont={bf,color=techpurple},labelsep=colon}
\lstset{basicstyle=\ttfamily\small,breaklines=true,frame=single,rulecolor=\color{techline},backgroundcolor=\color{techlavender},showstringspaces=false,keywordstyle=\color{techpurple}\bfseries,commentstyle=\color{gray}}
\newtcolorbox{techbox}[1][]{colback=techlavender,colframe=techline,arc=3mm,boxrule=.6pt,left=7pt,right=7pt,top=6pt,bottom=6pt,#1}
\newtcolorbox{warningbox}{colback=techpink!7,colframe=techpink,arc=3mm,boxrule=.6pt,left=7pt,right=7pt,top=6pt,bottom=6pt}
\newcommand{\TechFigure}[4]{%
\begin{figure}[H]\centering
\includegraphics[width=\linewidth,height=#2,keepaspectratio]{#1}
\caption{#3}\label{#4}
\end{figure}}
\newcommand{\TechCover}[3]{%
\begin{titlepage}
\begin{tikzpicture}[remember picture,overlay]
\fill[techlavender] (current page.south west) rectangle (current page.north east);
\fill[techpurple] (current page.north west) rectangle ([yshift=-3.8cm]current page.north east);
\fill[techpurple2] ([xshift=-3.4cm]current page.south east) circle (4.4cm);
\draw[techpurple2,line width=2pt] ([xshift=1.5cm,yshift=-7.4cm]current page.north west) circle (2.2cm);
\draw[white,line width=1.5pt] ([xshift=1.5cm,yshift=-7.4cm]current page.north west) circle (1.45cm);
\end{tikzpicture}
\vspace*{0.65cm}\color{white}\raggedright
{\fontsize{24}{28}\selectfont\bfseries TECH-FLOW\par}
{\large Sistema de gestión de inventario y ventas\par}
\vspace{3.2cm}\color{techink}
{\fontsize{30}{35}\selectfont\bfseries #1\par}
\vspace{0.35cm}{\Large\color{techpurple}#2\par}
\vfill
\begin{techbox}[colback=white,colframe=techline]
\textbf{Universidad:} Universidad Mayor de San Simón\\
\textbf{Docente:} CORINA JUSTINA FLORES VILLARROEL\\
\textbf{Integrantes:}\\
Jorge Henry Quispe Escalera\\
Logan Betho Canelas Silvestre\\
Fausto Jafeth Vilches Mendieta\\
Eunice Melany Zurita Pérez\\
\textbf{Entrega:} 30 de noviembre de 2026\\
\textbf{Defensas:} 1, 2 y 8 de diciembre de 2026\\
\end{techbox}
\end{titlepage}}

\begin{document}
\TechCover{Manual de usuario}{Cómo operar Tech-Flow: administrador, cajero y almacén}{Manual de usuario}
\tableofcontents\clearpage

\section{Conozca Tech-Flow}
Tech-Flow es el sistema que acompaña su jornada en la tienda: punto de venta y garantías, control de inventario y entradas, y la supervisión administrativa del día. Todo ocurre en una única ventana; los mensajes, formularios y confirmaciones se muestran dentro de ella. Puede alternar entre tema claro y oscuro con el botón de la pantalla de acceso; el tema elegido se conserva mientras la sesión siga abierta.

\section{Inicio de sesión}
\TechFigure{../../documentacion/capturas/01-login.png}{8.2cm}{Pantalla de inicio de sesión, selección de rol y cambio de tema.}{fig:user-login}

Seleccione su rol (Administrador, Cajero o Almacén), escriba su nombre de usuario y su contraseña y pulse el botón de inicio de sesión. Si alguno de los datos no corresponde, se mostrará el mensaje correspondiente; corrija y vuelva a intentarlo. Al entrar, cada rol verá su propio panel.

\section{Rol Administrador}

\subsection{Gestión de caja}
\TechFigure{../../documentacion/capturas/02-admin-caja.png}{8.2cm}{Apertura de caja desde el panel de administrador.}{fig:user-caja}

Al entrar, se propone abrir la caja. Si ya está abierta, se informa del estado actual. Sin caja abierta ningún cajero puede confirmar ventas; abra la caja al iniciar la jornada y ciérrela al final.

\subsection{Ventas y reportes}
\TechFigure{../../documentacion/capturas/03-admin-ventas.png}{8.2cm}{Panel de ventas del administrador y acción de exportación.}{fig:user-admin-ventas}

La pestaña \emph{Ventas} muestra las ventas confirmadas del día con su total y un botón \emph{Exportar informe} que guarda un archivo .txt con el desglose. La pestaña \emph{Entradas} hace lo mismo con las entradas recibidas y el costo de compra acumulado.

\subsection{Inventario}
La pestaña \emph{Inventario} ofrece dos vistas seleccionables: \emph{Inventario completo} (tabla con todos los productos; filtre por categoría con el ComboBox) y \emph{Productos por agotarse} (cinco unidades o menos, con contador de casos críticos).

\subsection{Empleados, productos y proveedores}
\begin{itemize}
\item \textbf{Empleados}: alta con identificador, nombre, rol y contraseña; baja del seleccionado. No se puede eliminar al último administrador ni el sistema permite duplicar el CI ni dos empleados con el mismo nombre.
\item \textbf{Productos}: alta con identificador generado automáticamente y edición de precios y stock; baja del seleccionado. Complete nombre, marca, categoría, descripción, precios y stock.
\item \textbf{Proveedores}: alta con nombre y teléfono de contacto y baja del seleccionado.
\end{itemize}

\section{Rol Cajero}

\subsection{Carrito y venta}
\TechFigure{../../documentacion/capturas/04-cajero-carrito.png}{8.2cm}{Carrito de venta con producto seleccionado.}{fig:user-carrito}

\begin{enumerate}
\item Escriba el código, nombre o marca del producto en el buscador y pulse \emph{BUSCAR}. Opcionalmente, refine por categoría con el ComboBox \emph{Categoría}.
\item Ajuste la cantidad y pulse ``+'' para añadir al carrito; el sistema descuenta visualmente el stock disponible.
\item En \emph{Carrito} revise las líneas, quite lo que no corresponda o vacíe con \emph{Vaciar carrito}.
\item Introduzca nombre, CI y teléfono del comprador (los tres son obligatorios).
\item Pulse \emph{Confirmar venta}; o \emph{Hacer proforma} para guardar una cotización .txt sin confirmar ni descontar stock.
\end{enumerate}

Tras confirmar, se muestra un cuadro con el resumen (se puede guardar el comprobante .txt o cerrar). Para anular una compra o proforma, ingrese el identificador y confirme: las unidades vuelven a existir si estaba confirmada.

\subsection{Garantías}
Pulse \emph{Consultar garantía} e introduzca el código de garantía, la venta o la orden: se indica si el equipo conserva cobertura y hasta cuándo.

\subsection{Importar una orden}
Seleccione un CSV/TXT con columnas \emph{nombreCliente}, \emph{idProducto}, \emph{cantidad}, \emph{ciCliente}, \emph{telefonoCliente} (\emph{idOrden} opcional). Se muestra una vista previa de las filas antes de confirmar la importación; nada se persiste hasta que confirma.

\section{Rol Almacén}

\subsection{Consultar inventario}
\TechFigure{../../documentacion/capturas/05-almacen-inventario.png}{8.2cm}{Inventario del panel de almacén.}{fig:user-inventario}

La tabla muestra ID, nombre, marca, categoría, precios y existencias; las existencias bajas se resaltan en rojo y las inválidas en gris. Filtre por categoría con el ComboBox de la parte superior derecha.

\subsection{Registrar entradas}
\TechFigure{../../documentacion/capturas/06-almacen-entradas.png}{8.2cm}{Búsqueda y registro de entradas y stock.}{fig:user-entradas}

Abra \emph{ENTRADAS Y STOCK}, busque el producto y añádalo con la cantidad recibida; confirme para registrar la entrada y aumentar el inventario. También puede importar un CSV de entradas con vista previa.

\subsection{Registrar producto}
Complete ID (generado), nombre, marca, categoría, descripción, precios y stock inicial. El precio de venta debe ser mayor que el de compra.

\section{Mensajes y solución de problemas}
\begin{itemize}
\item \textbf{``El carrito está vacío''}: añada productos antes de confirmar.
\item \textbf{``Stock insuficiente''}: no hay unidades suficientes; ajuste la cantidad.
\item \textbf{``No se pudo preparar la venta''}: datos del comprador o stock revisables; el cuadro indica qué falta.
\item \textbf{``La caja está cerrada''}: un administrador debe abrirla primero.
\item \textbf{Errores de archivo}: el destino del .txt no es escribible o la extensión no termina en \texttt{.txt}; guarde en otra carpeta.
\end{itemize}

\section{Cerrar sesión y recomendaciones}
Use \emph{Cerrar sesión} en la barra lateral para volver al acceso. Al terminar la jornada, el administrador cierra la caja y exporta los informes del día. Nunca comparta su contraseña: el sistema no recupera accesos.
\end{document}
